\section{Evaluation of the reference tree}\label{sec:reference}

We now evaluate the boundary contribution left open by the continuous
benchmark in Section~\ref{sec:functions}.  Above exponent $2$, the
reference tree follows the standard admission rule.  At cutoff $2$,
its remaining finite polynomial differs from the linear-sieve benchmark.
The occupation estimates of Section~\ref{sec:prime-paths} transport this
boundary discrepancy to the root, where it contributes at scale $B^{-2}$.

The sign requires a separate calculation.  We first compare with the
full product over the boundary exponents, whose signed contribution can
be evaluated exactly.  The stronger admission rule below exponent $2$
omits lower children and reduces this contribution.  Expanding at the
first omitted child will bound the loss and leave enough positive mass
to overcome the root's negative benchmark.  The resulting lower bound
is Proposition~\ref{prop:reference-margin}.  We also obtain a local
reference estimate for the later stopped-subtree comparison.

We begin by replacing the continuous factor $1/b$ in the benchmark by
the corresponding prime product and controlling the resulting quadrature
error.  After this error is removed, only the cutoff-$2$ discrepancy
remains to be evaluated.

For an available cutoff $b$, let $V(b)$ be the product of $1-1/p$ over
the available primes greater than $w$.  The root cutoff includes its
upper endpoint, whereas a cutoff at a previously chosen prime excludes
that prime.  Write $V_x$ for the latter product when the new prime has
exponent $x$.  The auxiliary cutoff $2$ includes the boundary primes
$w<p\leq w^2$.  Mertens' theorem and Lemma~\ref{lem:prime-harmonic} give
\begin{equation}\label{eq:ref-product}
 V(b)=\frac{C_w}{b}\bigl(1+O(E(b))\bigr),\qquad
 C_w=\frac{e^{-\gamma}}{\log w\,V_0}\longrightarrow1
 \quad(b\geq2).
\end{equation}
The formula is uniform in these endpoint conventions: excluding an
endpoint prime changes the product by $1+O(1/p)$, an error absorbed
by $E(b)$.

With $f_\ee=f$ and $f_\oo=F$ as before, benchmark a state by
$V(b)f_i(r/b)$.  The extension of $f$ below $2$ is used only at a
starting lower state, as in Section~\ref{sec:functions}.  An omitted
lower child always has benchmark zero, including when its formal ratio
is below the domain of the ordinary sieve functions.

\subsection{The quadrature error above the boundary}

Define $g_i=f_i-1$, with the starting extension just described, and set
$g_\ee(t)=-1$ whenever the prospective lower child is omitted.  At a
state with cutoff $b>2$, subtracting constants by the identity
\[
 V(b)=V(2)-\sum_{2<x(p)\leq b}\frac{V_{x(p)}}p
\]
(with the actual upper endpoint convention) gives the residual
\begin{equation}\label{eq:ref-residual}
 R_i^{(w)}(r,b)=V(2)g_i(r/2)-V(b)g_i(r/b)
       -\sum_{2<x(p)\leq b}\frac{V_{x(p)}}p
                          g_{i'}(r/x(p)-1).
\end{equation}
The sum here includes the omitted children with their deviation $-1$;
they are not included in any subsequent path propagation.  This
convention accounts exactly for the constants of the omitted terms.

The continuous analogue of \eqref{eq:ref-residual}, obtained by
replacing $V$ by $C_w$ divided by its exponent and the prime measure
by $dx/x$, is zero.  Subtracting the constant part from
\eqref{eq:functions-benchmark-cutoff}, with the same starting extension
and zero benchmark at omitted children, gives
\begin{equation}\label{eq:ref-differentiation}
 \frac{d}{db}\left(\frac{g_i(r/b)}b\right)
       =-\frac{g_{i'}(r/b-1)}{b^2}.
\end{equation}
Integrating this identity from $2$ to $b$ proves the cancellation.

\begin{lemma}[Summed quadrature residual]\label{lem:ref-residual}
Along the standard prime paths with $x>2$ from the root,
\begin{equation}\label{eq:ref-residual-sum}
 \sum_{\mathfrak p}h(\mathfrak p)
       \left|R_{i(\mathfrak p)}^{(w)}
                   (r(\mathfrak p),b(\mathfrak p))\right|=o(B^{-2}).
\end{equation}
The same assertion is uniform for lower starting ratios in
$[1.99,2.3]$ and starting cutoffs $B$ tending to infinity with the
parameters of \eqref{eq:parameters}.
\end{lemma}

\begin{proof}
The absolute next deviation in \eqref{eq:ref-residual} is nondecreasing
as a function of $x$.  Thus partial summation on an exponent interval
with lower endpoint $u\geq2$ bounds the error in replacing
\[
 \sum\frac1p\frac{g_{i'}(r/x(p)-1)}{x(p)}
 \quad\hbox{by}\quad
 \int\frac{g_{i'}(r/x-1)}{x^2}\,dx
\]
by $O(E(u)D/u)$, where $D$ is the maximum absolute next deviation on
that interval.  To justify the variation bound despite the additional
factor $1/x$, use the product-variation inequality: the total variation
is at most a constant times $D/u$.  By \eqref{eq:ref-product}, replacing
the products introduces endpoint errors of the same kind and a sum
error at most $O((\log B)E(u)D)$.  These estimates apply to subintervals
as well, and their continuous contributions cancel by
\eqref{eq:ref-differentiation}.

At a state of ratio $s$, all the absolute parent, boundary, and next
deviations under discussion are $O(\phi_i(s))$.  For example,
\[
 a(s)\leq\phi_\ee(s)/s,\qquad
 c(s-1)\leq\phi_\ee(s)/s,
\]
and the analogous inequalities follow from
$\phi_\oo(s)=s(c(s)+a(s-1))$ on its normal range.  Monotonicity bounds
the later arguments by these earliest ones.  The bounded extensions
at the initial states obey the same inequalities with an absolute
constant.  An omitted lower deviation is just $1$ in absolute value,
which occurs where the relevant upper derivative weight is bounded
away from zero.

First restrict to histories with all ratios at most
$S=(\log B)^2$.  If $b\leq(\log B)^4$, then $r\leq(\log B)^6$, and
the preceding partial summation bounds imply, for some fixed $C$,
\begin{equation}\label{eq:ref-small-cutoff-error}
 \frac{r^2|R_i^{(w)}(r,b)|}{\phi_i(s)}
       \ll(\log B)^C E(2).
\end{equation}
This is smaller than every negative power of $\log B$.

If $b>(\log B)^4$, split the exponent interval at $\sqrt b$.  Above
this point $E(\sqrt b)$ is smaller than every $B^{-A}$: indeed
$\sqrt{\sqrt b\log w}\gg(\log B)^{3/2}$.  Below it every relevant
next ratio is at least $s\sqrt b-1$, and the boundary ratio is at
least $s\sqrt b$.  The Dickman bounds from
Lemma~\ref{lem:derivative-weights}, together with the backward-ratio
bounds there, show that these deviations divided by $\phi_i(s)$
are smaller than every $B^{-A}$, uniformly for $s\leq S$.  One can see
the uniformity by comparing the arguments to $s$: their separation
is at least a constant times $(\log B)^2$, while the derivative weights
have exponentially decreasing ratios beyond any fixed compact set.
Any fixed polynomial losses are absorbed.  Since $r=O(B)$ along a
path, the same assertion holds after multiplication by $r^2$.

Apply \eqref{eq:prime-weight} and sum over at most
$N=O((\log B)^3)$ steps.  The starting factor
$\phi_\ee(s_{\mathrm{st}})/r_{\mathrm{st}}^2$ is $O(B^{-2})$, and
\eqref{eq:ref-small-cutoff-error} and its large-cutoff counterpart show
that the low-state contribution to \eqref{eq:ref-residual-sum} is
$o(B^{-2})$.

It remains to bound histories having a ratio above $S$.  The crude
bound for every residual is $O(\log B)$, since all deviations in its
definition are bounded and the total harmonic prime sum is
$O(\log B)$.  We claim the sharper bound
\begin{equation}\label{eq:ref-high-gap-error}
 |R_i^{(w)}(r,b)|=O_A(B^{-A})\qquad(r\geq S/2)
\end{equation}
without a ratio restriction.  In the preceding error estimates split
instead at
$Q=(\log B)\sqrt{\log\log B}$.  For $x\geq Q$, the PNT saving has
exponent at least a constant times
$(\log B)(\log\log B)^{1/4}$, and hence beats every $B^{-A}$.  For
$x\leq Q$, each relevant argument $r/x$ or $r/x-1$ is at least a
constant times $\log B/\sqrt{\log\log B}$.  The Dickman bound
$\rho(v)\leq\exp(-v\log v+O(v))$ makes its deviation smaller than every
$B^{-A}$.  The endpoint errors fall into the same alternatives.
This proves \eqref{eq:ref-high-gap-error}.

The total prefix weight is $O(B)$, so
\eqref{eq:ref-high-gap-error} deals with the high-history prefixes
whose terminal gaps are at least $S/2$.  A high ancestor on the
$x>2$ tree has gap greater than $2S$.  Any remaining prefix with
terminal gap below $S/2$ has reduced that ancestor's gap by more than
half.  The factorial-tail argument in Lemma~\ref{lem:high-states},
multiplied by the crude $O(\log B)$ residual bound, deals with those
prefixes.  This completes the proof.
\end{proof}

\subsection{The boundary anomaly}

For the actual discrete reference tree at cutoff $2$, put
\begin{equation}\label{eq:ref-anomaly}
 \delta_i^{(w)}(r)=P_i(r,2)-V(2)f_i(r/2),
 \qquad r\geq4\ (i=\ee),\quad r\geq2\ (i=\oo).
\end{equation}
These domains contain every state reached when iterating the root
recurrence only through exponents $x>2$.  A possible first upper state
arising from a starting ratio below $2$ still has gap tending to
infinity, so it causes no exception.

Let $Q_i(r,b)$ denote the continuous reference polynomial, obtained
from the same admission rule by replacing prime sums by ordered
integrals in $dx/x$, with exponents above $1$.  On the displayed
domains define
\begin{equation}\label{eq:ref-cont-anomaly}
 \delta_i(r)=Q_i(r,2)-\tfrac12 f_i(r/2).
\end{equation}

\begin{lemma}[Boundary convergence and decay]\label{lem:ref-anomaly}
On every fixed compact subset of their respective domains,
$\delta_i^{(w)}$ converges uniformly to the continuous function
$\delta_i$.  Uniformly for sufficiently large $w$,
\begin{equation}\label{eq:ref-anomaly-decay}
 |\delta_i^{(w)}(r)|+|\delta_i(r)|
       \ll\exp(-c r\log(r+2)).
\end{equation}
\end{lemma}

\begin{proof}
On a fixed compact gap interval the admitted polynomial has bounded
length, because every exponent exceeds $1$.  For each such length,
Lemma~\ref{lem:prime-harmonic} gives convergence of the ordered prime
sum to its harmonic integral.  Ties, moving upper endpoints, and
admission equalities are confined to sets of continuous measure zero.
The same argument works along any convergent sequence of gap
parameters.  The continuous integrals depend continuously on those
parameters by dominated convergence, so compactness yields uniform
convergence.  Also $V(2)\to1/2$ by \eqref{eq:ref-product}.

For decay, compare the cutoff polynomial with the full product $V(2)$.
At a first omitted lower leaf after $m$ boundary exponents, all of
which are at most $2$, the failure of admission requires
$m>(r-4)/2$.  The total discrepancy is bounded by the sum of the
weights of such first omissions times their exact subsequent survivor
product, which is at most $1$.  Thus it is bounded by
\[
 \sum_{m>(r-4)/2}\frac{T_2^m}{m!},\qquad
 T_2=\sum_{w<p\leq w^2}\frac1p=\log2+o(1).
\]
This is $O(\exp(-c r\log(r+2)))$.  The continuous full product is
$\exp(-\int_1^2 dx/x)=1/2$, and its first-omission decomposition has
the same bound.  Finally the benchmark deviations $a(r/2)$ and
$c(r/2)$ have the required decay by
Lemma~\ref{lem:future-integrals}.  Combining these comparisons proves
\eqref{eq:ref-anomaly-decay}.
\end{proof}

Iterate the difference between the reference recurrence and its
benchmark only through $x>2$.  At every visited state, including the
root, insert its anomaly \eqref{eq:ref-anomaly} and its residual
\eqref{eq:ref-residual}; propagate an insertion with the original
prefix weight and sign $(-1)^m$ at depth $m$.  More explicitly,
subtracting the benchmark from the recurrence gives its cutoff-$2$
anomaly minus the child differences, plus the residual; repeated
substitution gives
\begin{equation}\label{eq:ref-root-decomposition}
 P_\ee(r_0,B)=V(B)f_{\mathrm{ext}}(r_0/B)
                  +\mathcal D_w+o(B^{-2}),
 \qquad
 \mathcal D_w=\sum_{\mathfrak p}(-1)^{|\mathfrak p|}
                       h(\mathfrak p)\delta_{i(\mathfrak p)}^{(w)}
                                             (r(\mathfrak p)),
\end{equation}
where Lemma~\ref{lem:ref-residual} bounds the residual sum absolutely.

Compact convergence in Lemma~\ref{lem:ref-anomaly}, the compact
occupation limit in Proposition~\ref{prop:prime-occupation}, and its
exponentially weighted tail estimate justify passing to the limit in
$\mathcal D_w$.  The limiting stationary densities are
$\pi_\ee(t)=\phi_\ee(t)$ and
$\pi_\oo(t)=(1-t^{-2})\phi_\oo(t)$, and the arrived-at cutoff
restriction is $t<r/2$.  Moreover
\[
 \frac{B^2\phi_\ee(r_0/B)}{r_0^2}\longrightarrow
                         \frac{\phi_\ee(2)}4=e^\gamma.
\]
Consequently
\begin{align}\label{eq:ref-correction-limit}
 \lim_{w\to\infty}\frac{B^2\mathcal D_w}{e^\gamma}
  =\frac1{M_g}\biggl\{
   &\int_4^\infty r^2\delta_\ee(r)
                   \int_2^{r/2}dt\,\frac{dr}{r}\notag\\
   -&\int_2^\infty r^2\delta_\oo(r)
                   \int_1^{r/2}(1-t^{-2})\,dt\,\frac{dr}{r}
                   \biggr\}.
\end{align}

\subsection{The signed correction integral}

Set $y=r/2-1$ and define
\[
 l(y)=2Q_\ee(2y+2,2)\quad(y\geq1),\qquad
 u(y)=2Q_\oo(2y+2,2)\quad(y\geq0).
\]
The even inner integral in \eqref{eq:ref-correction-limit} is $y-1$;
the odd inner integral is $y^2/(y+1)$.  Since $r\,dr=4(y+1)\,dy$,
that limit is $2I/M_g$, where
\begin{equation}\label{eq:ref-I}
 I=\int_1^\infty(y^2-1)\bigl(l(y)-f(y+1)\bigr)\,dy
       -\int_0^\infty y^2\bigl(u(y)-F(y+1)\bigr)\,dy.
\end{equation}
All these integrals converge absolutely by
\eqref{eq:ref-anomaly-decay}.

To determine the sign of $I$, compare the boundary polynomials with
the full survivor product over exponents in $[1,2]$.  That product is
$1/2$, so this comparison replaces both $l$ and $u$ by $1$.
Denote the resulting integral by $I_{\mathrm{full}}$.  As $a=1$ up to
$2$ and $a+c=c_0$, it equals
\begin{equation}\label{eq:ref-Ifull}
 2e^\gamma\int_0^\infty\frac{y^2\rho(y)}{y+1}\,dy
           -\frac13-\int_2^\infty a(t)\,dt=\frac83.
\end{equation}
Here are the identities giving the last equality.  The Dickman
integral equation holds in the form
$y\rho(y)=\int_{\max(0,y-1)}^y\rho(t)\,dt$ also below $1$.
Integration and Fubini give
$\int_0^\infty y\rho(y)\,dy=\int_0^\infty\rho(y)\,dy$.
The differential equation gives
$\int_0^\infty\rho(y)/(y+1)\,dy=-\int_1^\infty\rho'(t)\,dt=1$.
Thus the decomposition
$y^2/(y+1)=y-1+1/(y+1)$ makes the first integral in
\eqref{eq:ref-Ifull} equal to $2e^\gamma$.  Finally
Lemma~\ref{lem:future-integrals}, at $s=1$ in the identity for $c$,
gives $\int_0^\infty a=2e^\gamma-1$; subtracting the interval $[0,2]$
gives $\int_2^\infty a=2e^\gamma-3$.

The actual boundary rule can only reduce this full-product contribution:
a first omitted lower leaf removes a positive term for a lower start
and a negative term for an upper start.  We quantify the resulting loss
by recording the first omission.  Consider ordered boundary exponents
\[
 2\geq x_1\geq\cdots\geq x_m\geq1,\qquad
 A_m=\sum_{j=1}^m x_j+x_m,\qquad
 A'=
 \begin{cases}
  0,&m=1,\\
  2,&m=2,\\
  A_{m-2},&m\geq3.
 \end{cases}
\]
For boundary exponents the stronger admission rule is
$r_{\mathrm{child}}\geq x+2$.  Starting at gap $2y+2$, this says
$2y\geq A_m$ at a prospective lower child.  The thresholds increase
within each parity: for $m\geq3$ their difference is
$x_{m-1}+2x_m-x_{m-2}>0$ on $[1,2]$, and $A_2>2$.
Consequently the first omitted child occurs at even $m$ for a lower
start, and at odd $m$ for an upper start, exactly when
\begin{equation}\label{eq:ref-omission-window}
                     A'/2\leq y<A_m/2.
\end{equation}
The exceptional value $A'=2$ for $m=2$ expresses the lower-start
domain $y\geq1$; $A'=0$ for $m=1$ expresses the upper-start domain
$y\geq0$.

After this first omitted leaf, the exact full survivor product over
smaller boundary exponents is $1/x_m$.  This follows either by the
product identity or by the absolutely convergent full inclusion--
exclusion series with harmonic intensity $\log x_m$.  Omitting a
positive even-depth term lowers the lower polynomial; omitting a
negative odd-depth term raises the upper polynomial.  These signs,
the factor $2$ in $l,u$, and integration over
\eqref{eq:ref-omission-window} give
\begin{equation}\label{eq:ref-loss-series}
 I_{\mathrm{full}}-I
 =2\sum_{m\geq1}\int_{2\geq x_1\geq\cdots\geq x_m\geq1}
    \frac1{x_m}
     \left\{\frac{A_m^3-(A')^3}{24}
             -\1_{2\mid m}\frac{A_m-A'}2\right\}
          \prod_{j=1}^m\frac{dx_j}{x_j}.
\end{equation}
Every term is nonnegative: its expression in braces is the integral
of $y^2$ or $y^2-1$ on its allowed window.  The bounds below also
establish convergence, so the first-omission summations and integrals
are justified by Tonelli's theorem.

Write $L_m$ for the $m$th term of \eqref{eq:ref-loss-series} and
$\lambda=\log2$.  The first term is
$L_1=\int_1^2(2x/3)\,dx=1$.  For completeness we evaluate the next
two terms by explicit one-dimensional reductions.  Integrating $x_1$
first in the $m=2$ term gives
\begin{align}\label{eq:ref-L2-reduction}
 L_2=\int_1^2\biggl\{
 &-\frac{23t}{18}+2+\frac2t-\frac{16}{9t^2}\notag\\
 &+\left(\frac{2t}3-\frac2t+\frac4{3t^2}\right)
                  \log\frac2t\biggr\}\,dt
       =-\frac{35}{36}-\lambda^2+3\lambda.
\end{align}
Indeed the original integrand in this case is
\[
 \frac1{x_1x_2^2}
 \left\{\frac{(x_1+2x_2)^3-8}{12}-(x_1+2x_2-2)\right\}.
\]
For $m=3$, the original integrand is
\[
 \frac{(x_1+x_2+2x_3)^3-8x_1^3}{12x_1x_2x_3^2}.
\]
Integrating $x_1$ from $x_2$ to $2$ and then $x_2$ from $x_3$ to $2$
yields
\begin{align}\label{eq:ref-L3-reduction}
 L_3=\int_1^2\biggl\{
 &\frac t3\log^2\frac2t
 +\left(-\frac{23t}{18}+2+\frac1t-\frac{14}{9t^2}\right)
                      \log\frac2t\notag\\
 &+\frac{119t}{216}-\frac94+\frac3{2t}+\frac{43}{27t^2}
                  \biggr\}\,dt
       =\frac{623}{432}-\frac{19}{12}\lambda+\frac13\lambda^2.
\end{align}
These equalities involve only elementary antiderivatives.  For an
explicit check of the last integrations, put
$J_{j,k}=\int_1^2 t^{j-1}\log^k(2/t)\,dt$.  The needed values are
\[
\begin{array}{c|rrrr}
 j&2&1&0&-1\\ \hline
 J_{j,0}&3/2&1&\lambda&1/2\\
 J_{j,1}&3/4-\lambda/2&1-\lambda&\lambda^2/2&\lambda-1/2
\end{array}
 \qquad J_{2,2}=\frac34-\frac\lambda2-\frac{\lambda^2}2.
\]
Substitution in \eqref{eq:ref-L2-reduction} and
\eqref{eq:ref-L3-reduction} proves the displayed formulas exactly.

For $m\geq4$, ordering gives
\[
 A_m-A'=x_{m-1}+2x_m-x_{m-2}\leq2x_m,\qquad
 A_m\leq(m+1)\overline x,\qquad
 \overline x=\frac1m\sum_{j=1}^m x_j.
\]
Since $A_m^3-(A')^3\leq3A_m^2(A_m-A')$, and dropping the negative
even term only increases the answer, the integrand including the
factor $2/x_m$ is at most $A_m^2/2$.  This last upper bound can be
replaced by the symmetric function $(m+1)^2\overline x^{\,2}/2$.
Under the probability measure $dx/(\lambda x)$ on $[1,2]$,
\[
 \E x=\lambda^{-1},\qquad \E x^2=\frac3{2\lambda},\qquad
 \E\overline x^{\,2}=
 \lambda^{-2}+\frac1m\left(\frac3{2\lambda}-\lambda^{-2}\right).
\]
Integration over the ordered region divides a symmetric integral by
$m!$.  We obtain the rigorous tail bound
\begin{equation}\label{eq:ref-loss-tail}
 \sum_{m\geq4}L_m\leq
 \frac12\sum_{m\geq4}(m+1)^2\frac{\lambda^m}{m!}
       \left\{\lambda^{-2}
          +\frac1m\left(\frac3{2\lambda}-\lambda^{-2}\right)\right\}
       <0.34.
\end{equation}
Here is a coarse numerical verification with strict rational margins.
From $0.69<\lambda<0.70$, the expression in braces for $m\geq4$ is at
most
$3/(4(0.69)^2)+3/(8(0.69))<2.14$.  Also
\[
 \sum_{m\geq4}(m+1)^2\frac{(0.70)^m}{m!}
  =e^{0.70}\bigl((0.70)^2+3(0.70)+1\bigr)
        -1-4(0.70)-\frac92(0.70)^2-\frac83(0.70)^3<0.313.
\]
For example, $e^{0.70}<2.014$, obtained by summing its Taylor series
and bounding the remaining positive tail geometrically, already gives
the last strict inequality.  Their product with $1/2$ is less than
$0.34$.  The exact expressions above also give $L_2<0.66$ and
$L_3<0.52$ on $0.69<\lambda<0.70$.  Therefore
\begin{equation}\label{eq:ref-I-positive}
 I>\frac83-1-0.66-0.52-0.34>0.14.
\end{equation}

\begin{proposition}[Positive reference margin]\label{prop:reference-margin}
At the root of the reference tree,
\begin{equation}\label{eq:ref-margin}
 \frac{B^2P_\ee(r_0,B)}{e^\gamma}
       \longrightarrow -a_\star+\frac{2I}{M_g}>0.02.
\end{equation}
In particular, there is an absolute $c_L>0$ such that
$B^2P_\ee(r_0,B)\geq c_L$ for all sufficiently large $z$.
\end{proposition}

\begin{proof}
The parameter definitions imply
\[
 r_0=2B-a_\star+2-\frac{2\log L}{\log w}+o(1)
                  =2B-a_\star+o(1).
\]
Since $f'_{\mathrm{ext}}(2)=e^\gamma$, \eqref{eq:ref-product} gives
\[
 \frac{B^2V(B)f_{\mathrm{ext}}(r_0/B)}{e^\gamma}
                           \longrightarrow-a_\star.
\]
Equations \eqref{eq:ref-root-decomposition},
\eqref{eq:ref-correction-limit}, and \eqref{eq:ref-I} give the limit
in \eqref{eq:ref-margin}.  Finally $M_g\leq4e^\gamma$ by
\eqref{eq:mg-bound}.  The elementary bound
$\gamma\leq\sum_{j=1}^6j^{-1}-\log6<0.66$ implies $e^\gamma<1.95$.
Together with $a_\star=0.01$ and \eqref{eq:ref-I-positive}, this yields
\[
 -a_\star+\frac{2I}{M_g}
       >-0.01+\frac{0.28}{4(1.95)}>0.02.
\]
The eventual positive constant $c_L$ follows.
\end{proof}

\begin{remark}[A sharper bound for the correction integral]
\label{rem:ref-sharper-margin}
The same analytic majorant gives $I>0.2229793730$ using rational
arithmetic.  To record the bound explicitly, put
\[
 b=\frac{693147180559946}{10^{15}},\qquad
 J(t)=\frac{517}{432}-\frac{17t}{12}+\frac{2t^2}{3},
\]
\[
 t_m(t)=\frac{(m+1)^2}{2m!}
 \left[\left(1-\frac1m\right)t^{m-2}
       +\frac{3}{2m}t^{m-1}\right].
\]
The expansion
$\log2=2\sum_{j\ge0}((2j+1)3^{2j+1})^{-1}$ shows that
$0.69<\lambda=\log2<b$: after $j=15$ the remainder is at most
$2/(33\cdot3^{33}(1-1/9))$.
The exact formulas for $L_2,L_3$ and the tail majorant
\eqref{eq:ref-loss-tail} give
$I\ge J(\lambda)-\sum_{m\ge4}t_m(\lambda)$.
Each $t_m$ is increasing, whereas $J$ is decreasing on $[0.69,b]$.
For $m\ge13$,
\[
 \frac{t_{m+1}(b)}{t_m(b)}
 \le b\frac{(m+2)^2}{(m+1)^3}
 \le\frac{225b}{2744}<1.
\]
Indeed the remaining factor in $t_m(b)$ is
$1+(3b/2-1)/m$, which decreases with $m$.
Summing the geometric remainder and evaluating the finite rational
expression therefore yields
\[
 I>J(b)-\sum_{m=4}^{12}t_m(b)
          -\frac{t_{13}(b)}{1-225b/2744}
   >\frac{222979373}{10^9}.
\]
\end{remark}

\subsection{A uniform local comparison}

The stopped-node argument needs a different consequence of the same
boundary analysis.  Its ratios lie slightly above $2$, and its cutoffs
lie in a large but fixed window.  We need a uniform upper comparison
with $f(r/b)$ there, rather than the signed root limit.  The following
estimate permits the cutoff window to be chosen before the prime limit;
its $O(1/b)$ constant is independent of that window.

\begin{lemma}[Local reference estimate]\label{lem:local-reference}
For the continuous reference polynomial,
\begin{equation}\label{eq:ref-local-cont}
 Q_\ee(r,b)=\frac{f(r/b)}b+O(b^{-2})
 \qquad\left(b\longrightarrow\infty,\quad2\leq r/b\leq2.22\right),
\end{equation}
uniformly in the displayed ratios.  On every fixed cutoff window
$[b_0,b_1]$ with $b_0$ sufficiently large, the corresponding discrete
estimate is
\begin{equation}\label{eq:ref-local-prime}
 \frac{P_\ee(r,b)}{V(b)}
        =f(r/b)+O(1/b)+o_{w\to\infty}(1),
 \qquad b\in[b_0,b_1],\quad2\leq r/b\leq2.22.
\end{equation}
The constant in $O(1/b)$ is absolute, the last convergence is uniform
on the fixed window, and either prime endpoint convention is allowed.
\end{lemma}

\begin{proof}
In the continuous model the benchmark is $f(r/b)/b$ and the quadrature
residual is identically zero by \eqref{eq:ref-differentiation}.  Its
correction is the sum of the boundary anomalies $\delta_i$ along
the standard paths above $2$, with parity signs.  The continuous
version of \eqref{eq:path-weight} supplies the starting factor
$\phi_\ee(r/b)/r^2\ll b^{-2}$, uniformly for $r/b\in[2,2.22]$.
The expected sum of the absolute remaining reward
$R^2|\delta_i(R)|/\phi_i(s)$ is bounded uniformly in the start.
To verify this, use bounded compact-gap visits on a fixed compact
range and, outside it, sum the unit logarithmic gap bands from
Lemma~\ref{lem:band-bounds}.  After canceling $\phi_i$, a band
$R\leq r'\leq eR$ costs at most a fixed polynomial in $R$ times
$\exp(-cR\log(R+2))$, by \eqref{eq:ref-anomaly-decay}; the resulting
series converges.  The empty prefix and the explicitly bounded first
step satisfy the same estimate.  This proves
\eqref{eq:ref-local-cont}.

For fixed $b\in[b_0,b_1]$ and the indicated ratios, the starting gap
is bounded and the ordered reference polynomial has uniformly bounded
length.  Lemma~\ref{lem:prime-harmonic} therefore gives uniform
convergence of the discrete polynomial to $Q_\ee$: along any
convergent parameter sequence, tie, cutoff, and admission boundaries
have continuous measure zero, and dominated convergence applies to
each bounded-length integral.  Compactness makes this uniform on
the whole window.  Endpoint inclusion or exclusion makes no
difference to the limit.  Finally \eqref{eq:ref-product} gives
$V(b)\to1/b$ uniformly there.  Divide the continuous estimate by
$1/b$, and transfer it to the primes, to obtain
\eqref{eq:ref-local-prime}.
\end{proof}
